\subsection{Principal norm specialization}\label{sec:norms}

We recover primitive principal norm values by taking both forms in
the preceding results to be the norm form of one order. Thus the
ideal-class comparison and the local-density normalization require no
separate argument.

Fix an imaginary quadratic field $K$ of fundamental discriminant
$D<0$, and an order $\mathcal O=\Z+f\mathcal O_K$ of discriminant
$\Delta=f^2D$.  Choose a basis $1,\omega$ of $\mathcal O$ and put
\[
 Q(x,y)=\operatorname N_{K/\Q}(x+y\omega),\qquad
 P=\{Q(x,y)>0:x,y\in\Z,\ \gcd(x,y)=1\}.
\]
Primitivity refers to this order basis.  Write
\[
 \begin{gathered}
 \mathscr V_p=Q(\Z_p^2\setminus p\Z_p^2),\qquad
 \Lambda=\{n\ge1:n\in\mathscr V_p\text{ for every }p\},\\
 R_T(N)=\#\{(a,b)\in T^2:a+b=N\}.
 \end{gathered}
\]
Thus $P\subset\Lambda$, and $R_T$ counts numerical values without
coordinate multiplicities.  A target is locally admissible if
$N\in\mathscr V_p+\mathscr V_p$ for every prime $p$.

Let $S=\{p:p\mid2\Delta\}$, let $\chi=\chi_D$, and define
\[
 d_p(N)=\int_{\Z_p}\mathbf1_{\mathscr V_p}(a)
                    \mathbf1_{\mathscr V_p}(N-a)\,da\quad(p\in S),
 \qquad
 J_{\mathcal O}(N)=
 \prod_{\substack{p\mid N\\p\notin S,\ \chi(p)=-1}}
                         \frac{p-1}{p-2}.
\]
The Haar measure of $\Z_p$ is one.  The local densities concern the
set of values, rather than the number of their preimages under $Q$.

An independent weighted quadratic correlation and positive restoration
give an order-of-magnitude result. Its classical core estimate is proved
in full in Appendix~\ref{sec:older-route}, using no new zero-free or
weighted-dilation input.

\begin{corollary}[The quadratic order bounds]\label{thm:norm-order}
For every fixed imaginary quadratic order there are constants
$c_{\mathcal O},C_{\mathcal O}>0$ and $N_{\mathcal O}$ such that every
locally admissible $N\ge N_{\mathcal O}$ satisfies
\[
 c_{\mathcal O}\frac{NJ_{\mathcal O}(N)}{\log N}
       \le R_P(N)\le R_\Lambda(N)
       \le C_{\mathcal O}\frac{NJ_{\mathcal O}(N)}{\log N}.
\]
Every nonadmissible target has count zero.  These bounds use
Lemma~\ref{lem:older-quadratic-filter} and
Proposition~\ref{prop:norm-class-error}; they do not use
Theorem~\ref{thm:general}.
\end{corollary}

The leading-term theorem gives the following explicit specialization.

\begin{corollary}\label{thm:norm-leading}
Under the analytic inputs of Section~\ref{sec:inputs}, for every fixed
imaginary quadratic order, uniformly as $N$ tends to infinity through
the locally admissible targets,
\[
 R_P(N)\sim\mathfrak A_{\mathcal O}(N)
                     J_{\mathcal O}(N)\frac N{\log N},
\]
where the explicit positive coefficient is
\begin{equation}\label{eq:norm-coefficient}
 \mathfrak A_{\mathcal O}(N)=
 \frac{4L(1,\chi)}{\pi}
 \prod_{p\in S}\frac{1-\chi(p)/p}{1-1/p}
 \prod_{\substack{p\notin S\\\chi(p)=-1}}
                  \left(1-\frac2{p(p-1)}\right)
 \prod_{p\in S}d_p(N).
\end{equation}
Every nonadmissible target has count zero.  Each $d_p(N)$ is a
computable rational periodic function, and the infinite product
converges absolutely.  No restriction on the conductor or the ring
class group is required; the order remains fixed.
\end{corollary}

\begin{proof}[Specialization]
The norm form is primitive of discriminant \(\Delta\) and corresponds
to \(A=\mathcal O\), hence \(c_Q=0\). Lemma~\ref{lem:form-ideal-shadow}
therefore reads
\begin{equation}\label{eq:norm-shadow-ideals}
 n\in\Lambda\quad\Longleftrightarrow\quad
 \text{some primitive integral invertible }I
       \text{ has }n(I)=n,\ [I]\in2G,
\end{equation}
and replacing \(2G\) by zero characterizes \(P\).
The classical bounds just proved for mixed forms give
Corollary~\ref{thm:norm-order}. The two local filters now coincide,
\(d_{p;1,2}=d_p\), and \(J_\Delta=J_{\mathcal O}\), so
Theorem~\ref{thm:form-binary-leading} gives
\eqref{eq:norm-coefficient} and Corollary~\ref{thm:norm-leading}.
In particular, Proposition~\ref{prop:norm-class-error} gives
\(R_\Lambda-R_P\ll_{\mathcal O}NJ_{\mathcal O}/(\log N)^{6/5}\);
if \(2G=0\), it gives \(P=\Lambda\) exactly.
The finite local computation, prefix upper bound, summability and
Haar normalization are given by
\eqref{eq:norm-local-density-finite},
\eqref{eq:norm-prefix-tail} and \eqref{eq:form-haar-strata}.
They include all conductor-local factors in the principal case.
\end{proof}

\begin{remark}[Gaussian normalization]
For $\mathcal O=\Z[i]$ one has $P=C_4\sqcup2C_4$ and
\[
 d_2(N)=
 \begin{cases}
 1/4,&N\bmod8\in\{2,3,6,7\},\\
 1/8,&N\bmod8=4,\\
 0,&N\bmod8\in\{0,1,5\}.
 \end{cases}
\]
Consequently the coefficient of $J_{\mathcal O}(N)N/\log N$ is
$\frac12\prod_{p\equiv3\ (4)}(1-2/(p(p-1)))$ in the first four
classes and half that coefficient in the class $4\bmod8$.
This also calibrates the value-density normalization: the two Gamma
factors are already included in $c_{\mathcal C}^2$, and no additional
beta-integral factor occurs.
\end{remark}
